% GENERATED FILE. Edit canonical lessons, then run npm run generate:books.
% canonical-source-hash: fnv1a64:8f2cd09c0aa2c813
% canonical-lessons: TA-C93-meaning, TA-C93-intamil, TA-C93-howsay, TA-C93-show, TA-C93-didnthear

\chapter{How Do You Say It?}
\label{ch:ta-how-do-you-say-it}
\hlchaptermodality{93}

\begin{chapteropening}
\textbf{By the end of this chapter:} \emph{I can ask what something means and how to say it in Tamil, ask to be shown, and say I did not hear.}

Complete the last lesson of chapter 93: I can ask what something means and how to say it in Tamil, ask to be shown, and say I did not hear.
\end{chapteropening}

\section[\texorpdfstring{\ta{அர்த்தம்} (arttam)}{அர்த்தம் (arttam)}]{\ta{அர்த்தம்} (arttam) — meaning}

\label{lesson:TA-C93-meaning}



\begin{quote}
Before the new one: say the Tamil for one minute, then the Tamil for please wait.
\end{quote}

\subsection*{You'll want to know: \ta{அர்த்தம்}}
\textbf{\ta{அர்த்தம்}} (\emph{arttam}) — \textquotedblleft{}meaning\textquotedblright{}.

\textbf{\ta{இதன்} \ta{அர்த்தம்} \ta{என்ன}?} — \textquotedblleft{}what does this mean?\textquotedblright{}, the question that turns anyone into a teacher.

\textbf{In the family, and next door}

\noindent
\begin{tabularx}{\linewidth}{@{}>{\raggedright\arraybackslash}X>{\raggedright\arraybackslash}X>{\raggedright\arraybackslash}X@{}}
\toprule
\textbf{} & \textbf{word} & \textbf{same root} \\
\midrule
Kannada & \ka{ಅರ್ಥ} (\emph{artha}) & yes \\
Malayalam & \ml{അർത്ഥം} (\emph{arthaṁ}) & yes \\
Hindi (neighbour) & \dv{मतलब} (\emph{matlab}) &  \\
English & meaning &  \\
\bottomrule
\end{tabularx}

\begin{grammarlens}[title={What you've built}]
The question that unlocks a word.
\end{grammarlens}

\subsection*{Your turn}
\begin{itemize}
\raggedright
  \item \emph{Say it:} \emph{arttam}
  \item \emph{Say it:} \emph{arttam}, once more
  \item \emph{Say it:} ask \emph{itaṉ arttam eṉṉa?}
  \item \emph{Recall:} say the Tamil for one minute, then the Tamil for please wait, then say \emph{arttam} again
\end{itemize}

\begin{tcolorbox}[breakable,colback=teal!4,colframe=teal!35!black,title={Before you move on}]
What does \ta{அர்த்தம்} mean? (\textquotedblleft{}Meaning\textquotedblright{}.) Say it once more.
\end{tcolorbox}
\section[\texorpdfstring{\ta{தமிழில்} (tamiḻil)}{தமிழில் (tamiḻil)}]{\ta{தமிழில்} (tamiḻil) — in Tamil}

\label{lesson:TA-C93-intamil}



\begin{quote}
Before the new one: say the Tamil for please wait, then the Tamil for meaning.
\end{quote}

\subsection*{You'll want to know: \ta{தமிழில்}}
\textbf{\ta{தமிழில்}} (\emph{tamiḻil}) — \textquotedblleft{}in Tamil\textquotedblright{}.

\textbf{-\ta{இல்}} is \textquotedblleft{}in\textquotedblright{}, added to the name of the language. \textbf{\ta{ஆங்கிலத்தில்}} would be \textquotedblleft{}in English\textquotedblright{}.

\begin{grammarlens}[title={What you've built}]
A way to name the language you want.
\end{grammarlens}

\subsection*{Your turn}
\begin{itemize}
\raggedright
  \item \emph{Say it:} \emph{tamiḻil}
  \item \emph{Say it:} \emph{tamiḻil}, once more
  \item \emph{Say it:} say \emph{tamiḻil}
  \item \emph{Recall:} say the Tamil for please wait, then the Tamil for meaning, then say \emph{tamiḻil} again
\end{itemize}

\begin{tcolorbox}[breakable,colback=teal!4,colframe=teal!35!black,title={Before you move on}]
What does \ta{தமிழில்} mean? (\textquotedblleft{}In Tamil\textquotedblright{}.) Say it once more.
\end{tcolorbox}
\section[\texorpdfstring{\ta{எப்படிச்} \ta{சொல்வது}? (eppaṭic colvatu?)}{எப்படிச் சொல்வது? (eppaṭic colvatu?)}]{\ta{எப்படிச்} \ta{சொல்வது}? (eppaṭic colvatu?) — how do you say it?}

\label{lesson:TA-C93-howsay}



\begin{quote}
Before the new one: say the Tamil for meaning, then the Tamil for in Tamil.
\end{quote}

\subsection*{You'll want to know: \ta{எப்படிச்} \ta{சொல்வது}?}
\textbf{\ta{எப்படிச்} \ta{சொல்வது}?} (\emph{eppaṭic colvatu?}) — \textquotedblleft{}how does one say it?\textquotedblright{}.

\textbf{\ta{தமிழில்} \ta{இதை} \ta{எப்படிச்} \ta{சொல்வது}?} — \textquotedblleft{}how do you say this in Tamil?\textquotedblright{}. \textbf{\ta{எப்படி}} is the \textquotedblleft{}how\textquotedblright{} you know.

\textbf{In the family, and next door}

\noindent
\begin{tabularx}{\linewidth}{@{}>{\raggedright\arraybackslash}X>{\raggedright\arraybackslash}X@{}}
\toprule
\textbf{} & \textbf{word} \\
\midrule
Telugu & \te{ఎలా} \te{అంటారు}? (\emph{elā aṇṭāru?}) \\
Hindi (neighbour) & \dv{कैसे} \dv{कहते} \dv{हैं}? (\emph{kaise kahte haiṁ?}) \\
English & how do you say it? \\
\bottomrule
\end{tabularx}

\begin{grammarlens}[title={What you've built}]
The question that lets you learn from anyone.
\end{grammarlens}

\subsection*{Your turn}
\begin{itemize}
\raggedright
  \item \emph{Say it:} \emph{eppaṭic colvatu?}
  \item \emph{Say it:} \emph{eppaṭic colvatu?}, once more
  \item \emph{Say it:} ask \emph{tamiḻil itai eppaṭic colvatu?}
  \item \emph{Recall:} say the Tamil for meaning, then the Tamil for in Tamil, then say \emph{eppaṭic colvatu?} again
\end{itemize}

\begin{tcolorbox}[breakable,colback=teal!4,colframe=teal!35!black,title={Before you move on}]
What does \ta{எப்படிச்} \ta{சொல்வது}? mean? (\textquotedblleft{}How do you say it?\textquotedblright{}.) Say it once more.
\end{tcolorbox}
\section[\texorpdfstring{\ta{காட்டுங்கள்} (kāṭṭuṅkaḷ)}{காட்டுங்கள் (kāṭṭuṅkaḷ)}]{\ta{காட்டுங்கள்} (kāṭṭuṅkaḷ) — please show me}

\label{lesson:TA-C93-show}



\begin{quote}
Before the new one: say the Tamil for in Tamil, then the Tamil for how do you say it.
\end{quote}

\subsection*{You'll want to know: \ta{காட்டுங்கள்}}
\textbf{\ta{காட்டுங்கள்}} (\emph{kāṭṭuṅkaḷ}) — \textquotedblleft{}please show me\textquotedblright{}.

When a word will not come, point and say \textbf{\ta{காட்டுங்கள்}}.

\textbf{In the family, and next door}

\noindent
\begin{tabularx}{\linewidth}{@{}>{\raggedright\arraybackslash}X>{\raggedright\arraybackslash}X@{}}
\toprule
\textbf{} & \textbf{word} \\
\midrule
Kannada & \ka{ತೋರಿಸಿ} (\emph{tōrisi}) \\
Telugu & \te{చూపించండి} (\emph{cūpiñcaṇḍi}) \\
Hindi (neighbour) & \dv{दिखाइए} (\emph{dikhāie}) \\
English & please show me \\
\bottomrule
\end{tabularx}

\begin{grammarlens}[title={What you've built}]
A request for when words run out.
\end{grammarlens}

\subsection*{Your turn}
\begin{itemize}
\raggedright
  \item \emph{Say it:} \emph{kāṭṭuṅkaḷ}
  \item \emph{Say it:} \emph{kāṭṭuṅkaḷ}, once more
  \item \emph{Say it:} say \emph{kāṭṭuṅkaḷ}
  \item \emph{Recall:} say the Tamil for in Tamil, then the Tamil for how do you say it, then say \emph{kāṭṭuṅkaḷ} again
\end{itemize}

\begin{tcolorbox}[breakable,colback=teal!4,colframe=teal!35!black,title={Before you move on}]
What does \ta{காட்டுங்கள்} mean? (\textquotedblleft{}Please show me\textquotedblright{}.) Say it once more.
\end{tcolorbox}
\section[\texorpdfstring{\ta{கேட்கவில்லை} (kēṭkavillai)}{கேட்கவில்லை (kēṭkavillai)}]{\ta{கேட்கவில்லை} (kēṭkavillai) — I did not hear it}

\label{lesson:TA-C93-didnthear}



\begin{quote}
Before the new one: say the Tamil for how do you say it, then the Tamil for please show me.
\end{quote}

\subsection*{You'll want to know: \ta{கேட்கவில்லை}}
\textbf{\ta{கேட்கவில்லை}} (\emph{kēṭkavillai}) — \textquotedblleft{}I did not hear (it)\textquotedblright{}.

It ends in \textbf{\ta{இல்லை}}, the \textquotedblleft{}no\textquotedblright{} of your first chapter. Pair it with \textbf{\ta{சத்தமாக}}.

\textbf{In the family, and next door}

\noindent
\begin{tabularx}{\linewidth}{@{}>{\raggedright\arraybackslash}X>{\raggedright\arraybackslash}X@{}}
\toprule
\textbf{} & \textbf{word} \\
\midrule
Telugu & \te{వినిపించలేదు} (\emph{vinipiñcalēdu}) \\
Hindi (neighbour) & \dv{सुनाई} \dv{नहीं} \dv{दिया} (\emph{sunāī nahīṁ diyā}) \\
English & I did not hear it \\
\bottomrule
\end{tabularx}

\begin{grammarlens}[title={What you've built}]
Meaning, in Tamil, how to say it, show me, I did not hear.
\end{grammarlens}

\subsection*{Your turn}
\begin{itemize}
\raggedright
  \item \emph{Say it:} \emph{kēṭkavillai}
  \item \emph{Say it:} \emph{kēṭkavillai}, once more
  \item \emph{Say it:} say \emph{kēṭkavillai, cattamākap pēcuṅkaḷ}
  \item \emph{Recall:} say the Tamil for how do you say it, then the Tamil for please show me, then say \emph{kēṭkavillai} again
\end{itemize}

\begin{tcolorbox}[breakable,colback=teal!4,colframe=teal!35!black,title={Before you move on}]
What does \ta{கேட்கவில்லை} mean? (\textquotedblleft{}I did not hear it\textquotedblright{}.) Say it once more.
\end{tcolorbox}
